% ============================================================
%  Sage: Academic Paper — LaTeX Source
%  Template: generic two-column conference style
%  To switch venue:
%    ACM:  \documentclass[sigconf]{acmart}
%    IEEE: \documentclass[conference]{IEEEtran}
% ============================================================
\documentclass[10pt,twocolumn]{article}

\usepackage[margin=0.75in,top=1in,bottom=1in]{geometry}
\usepackage{times}
\usepackage{graphicx}
\usepackage{amsmath}
\usepackage{amssymb}
\usepackage{booktabs}
\usepackage{tabularx}
\usepackage{array}
\usepackage{multirow}
\usepackage{enumitem}
\usepackage{listings}
\usepackage{xcolor}
\usepackage{hyperref}
\usepackage{microtype}
\usepackage{caption}
\usepackage{abstract}
\usepackage{titlesec}
\usepackage{fancyhdr}
\usepackage{cite}

% Section formatting
\titleformat{\section}{\normalfont\large\bfseries}{\thesection.}{0.5em}{}
\titleformat{\subsection}{\normalfont\normalsize\bfseries}{\thesubsection}{0.5em}{}
\titlespacing*{\section}{0pt}{10pt}{4pt}
\titlespacing*{\subsection}{0pt}{8pt}{3pt}

% Code listing style
\lstset{
  basicstyle=\ttfamily\footnotesize,
  backgroundcolor=\color{gray!10},
  frame=single,
  framerule=0.4pt,
  rulecolor=\color{gray!50},
  breaklines=true,
  captionpos=b,
  numbers=none,
  keywordstyle=\color{blue!70!black},
  commentstyle=\color{green!50!black},
  stringstyle=\color{red!60!black},
}

% Compact list spacing
\setlist[itemize]{noitemsep, topsep=3pt, leftmargin=*}
\setlist[enumerate]{noitemsep, topsep=3pt, leftmargin=*}

\hypersetup{colorlinks=true, linkcolor=blue!60!black, citecolor=blue!60!black, urlcolor=blue!60!black}

% ============================================================
\begin{document}

\title{\textbf{Sage: A Reproducible, Shareable Natural Language\\
Science Agent for Interactive Data Exploration}}

\author{
[Author Names --- to be completed]\\
\small [Affiliation]\\
\small [Target Conference] $\bullet$ [Year]
}

\date{}
\maketitle

% ============================================================
\begin{abstract}
We present \textbf{Sage} (Science Agent), a natural language infrastructure for scientific data exploration that makes AI-driven analysis traceable, reproducible, and shareable. Deployable on any JupyterHub or local Jupyter environment and compatible with any major LLM provider, Sage allows researchers to express complete scientific workflows---from data discovery and integration through analysis and visualization---as plain English prompts, with no programming required. Every analysis step is logged in collapsible, inspectable tool panels embedded directly in the notebook, ensuring full auditability. A persistent conversational memory model enables coherent multi-step reasoning across notebook cells. A plug-in skill architecture integrated with the SkillsMP marketplace allows domain expertise to be acquired on-demand in natural language, without restarting the kernel. All outputs---GeoJSON maps, CSV tables, PNG charts---are saved to persistent, notebook-scoped directories, making results reproducible and shareable as standard \texttt{.ipynb} files. We demonstrate Sage on three real-world scientific use cases using data sources from the National Data Platform (NDP): (1) seismic event detection correlated with GNSS geodetic time series; (2) multi-dimensional flood impact assessment on the Kanawha River; and (3) on-the-fly emergency management skill installation and immediate use. Sage represents a new class of scientific tool: a recordable, shareable science agent---analogous to Claude Code or GitHub Copilot Workspace for software development, but designed for interactive scientific exploration with full auditability and reproducibility as first-class architectural properties.
\end{abstract}

% ============================================================
\section{Introduction}

The reproducibility crisis in computational science is well-documented. Studies estimating that a majority of published scientific results cannot be independently reproduced point to a common underlying cause: analyses are conducted in ad-hoc, undocumented workflows tightly coupled to a specific researcher's environment, personal code libraries, and tacit procedural knowledge. Even within the Jupyter notebook ecosystem---designed explicitly to combine narrative, code, and results---the gap between what is recorded and what is actually done remains wide. Cells are executed out of order, intermediate outputs are discarded, and environment dependencies go unspecified. The notebook has become a display medium rather than an audit trail.

Concurrently, the rise of AI coding agents---Claude Code, GitHub Copilot Workspace, Devin, and their kin---has demonstrated that large language models (LLMs) can dramatically accelerate software development by autonomously writing, executing, and debugging code. These tools are transformative for engineering but are not designed for science. They produce opaque artifacts: code that works, but whose provenance and reasoning are not captured in a form a colleague can inspect, validate, and reproduce. A software engineer using Claude Code delivers working software. A scientist using Claude Code delivers a script---not a reproducible experiment.

\textbf{Sage} (Science Agent) is a natural language infrastructure for scientific data exploration that runs on any standard Jupyter environment. Unlike AI coding assistants, which produce code as the deliverable, Sage treats generated code as a hidden implementation detail: the user sees only a plain English question, structured tool-call logs, and a final report with inline interactive maps, charts, and tables. The deliverable is a verifiable, shareable scientific record---not a script. Sage is LLM-agnostic, supporting OpenAI, Anthropic, and any OpenAI-compatible endpoint, and is deployable anywhere JupyterHub or JupyterLab can run---from a personal laptop to a national-scale research platform.

This paper makes the following contributions:
\begin{itemize}
  \item The \texttt{\%\%ask} / \texttt{\%reset} magic command architecture---a zero-configuration interface that turns any Jupyter notebook into a natural language science workbench.
  \item A Python-native cross-cell conversational memory model---enabling coherent multi-step scientific reasoning across notebook cells without a database checkpointer.
  \item A two-tier skill architecture---bundled domain skills at image level, plus runtime extensibility via the SkillsMP marketplace, all installable and usable in natural language within a single session.
  \item A notebook-scoped persistent output management system---with automatic file-type dispatch (GeoJSON, WMS, CSV, PNG) and multi-layer Folium map composition.
  \item Three real-world scientific demonstrations spanning geophysics, flood impact assessment, and emergency management, using publicly available data sources and a community skill marketplace.
\end{itemize}

The remainder of this paper is organized as follows. Section~2 surveys related work. Section~3 describes the system architecture. Section~4 presents the three demonstrations. Section~5 discusses design principles. Section~6 addresses limitations and future work. Section~7 concludes.

% ============================================================
\section{Related Work}

We organize related work into four categories: NL-to-code assistants in notebooks, autonomous notebook agents, domain-specific science agents, and skill/plugin marketplaces.

\subsection{Natural Language-to-Code Assistants in Notebooks}

The most prevalent class of AI notebook tools translates natural language into Python code that the user then reviews and executes. Jupyter AI---the official JupyterLab extension---supports multiple LLM backends and provides chat-based code generation and cell explanation, but delegates all execution to the user~\cite{jupyter_ai}. Microsoft Copilot in notebooks follows the same paradigm.

Sage takes a fundamentally different position: the generated code is an implementation detail, not a deliverable. The user sees only English input, structured tool-call logs, and English-plus-visual output. Code generation is a means, not an end.

\subsection{Autonomous Notebook Agents}

Several commercial platforms have moved toward more autonomous AI agents within notebooks. Deepnote's AI Agent autonomously creates SQL and Python blocks, executes them, inspects results, and self-corrects~\cite{deepnote}. Noteable's AI similarly generates and executes code within a notebook session. Google Colab has introduced agentic capabilities that can chain multiple code cells.

These systems share a common design: the agent's working medium is \emph{code cells}. Sage's working medium is \emph{natural language and structured visual output}. The distinction is consequential for reproducibility: a notebook full of AI-generated code cells is not more reproducible than a notebook full of human-written code cells, unless the reasoning behind each cell is also recorded.

\subsection{Domain-Specific Science Agents}

A growing body of work applies LLM agents to specific scientific domains. ChemCrow augments LLMs with chemistry tools for synthesis planning~\cite{chemcrow}. Bioinformatics agents have been built for single-cell RNA analysis. GeoAgent targets geospatial reasoning. These systems demonstrate domain viability but are domain-specific by design.

Sage is domain-agnostic at the system level; domain knowledge is encoded in skills. The same infrastructure that handles seismic analysis also handles flood impact assessment and emergency management---because skills, not system code, carry the domain logic.

\subsection{Skill and Plugin Marketplaces}

The concept of composable, reusable agent skills has gained traction with the SKILL.md open standard---a format for describing agent capabilities in natural language, subsequently adopted by the SkillsMP marketplace~\cite{skillsmp}, which currently hosts over 600,000 community-contributed skills. Sage integrates with SkillsMP as its primary extension mechanism, allowing users to discover, install, and immediately use marketplace skills within a live session.

\subsection{Summary}

Table~\ref{tab:comparison} summarizes the key distinguishing properties of Sage relative to the systems described above. The combination of NL-in/NL-out interface, traceable execution, persistent reproducible outputs, domain skill extensibility, scientific data API integration, standard \texttt{.ipynb} format, and deployment independence is not achieved by any existing system.

\begin{table*}[t]
\centering
\caption{Comparison of Sage with related AI notebook and agent systems. \checkmark~= fully supported, $\times$~= not supported, P~= limited or platform-dependent.}
\label{tab:comparison}
\small
\begin{tabular}{lcccccc}
\toprule
\textbf{Property} & \textbf{Sage} & \textbf{Jupyter AI} & \textbf{Colab Agent} & \textbf{Noteable} & \textbf{ChatGPT ADA} & \textbf{Deepnote} \\
\midrule
NL input/output (no code exposed)       & \checkmark & $\times$ & \checkmark & P & \checkmark & P \\
Steps traceable \& inspectable          & \checkmark & $\times$ & $\times$   & P & $\times$   & P \\
Persistent, reproducible outputs        & \checkmark & $\times$ & $\times$   & \checkmark & $\times$ & \checkmark \\
Domain skill plug-ins                   & \checkmark & $\times$ & $\times$   & $\times$ & $\times$   & $\times$ \\
Skills installable in natural language  & \checkmark & $\times$ & $\times$   & $\times$ & $\times$   & $\times$ \\
Scientific data APIs (WMS/WCS/WFS)      & \checkmark & $\times$ & $\times$   & $\times$ & $\times$   & $\times$ \\
Self-hostable (no cloud lock-in)        & \checkmark & $\times$ & $\times$   & $\times$ & $\times$   & $\times$ \\
Standard \texttt{.ipynb} format         & \checkmark & \checkmark & \checkmark & \checkmark & $\times$ & \checkmark \\
Cross-cell conversational memory        & \checkmark & $\times$ & $\times$   & P & $\times$   & \checkmark \\
\bottomrule
\end{tabular}
\end{table*}

% ============================================================
\section{System Architecture}

Sage is implemented as a single Python module---\texttt{sage\_magic.py}---that registers two IPython magic commands and orchestrates all agent execution, tool display, memory management, output handling, and file rendering. The module is installed as a JupyterHub startup script, requiring no per-notebook initialization.

\subsection{Deployment Stack}

The Sage Docker image is built on \texttt{quay.io/jupyter/base-notebook}, the standard Jupyter project base image providing Python 3, JupyterLab, and a pre-configured \texttt{jovyan} user environment. Scientific Python packages (GeoPandas, Folium, rasterio, matplotlib) are installed at build time. The image can be run locally via \texttt{docker run}, or deployed on any JupyterHub infrastructure. The demonstrations in this paper use the National Research Platform (NRP) JupyterHub, a Kubernetes-based multi-user environment, but this is not a requirement.

The underlying agent framework is \texttt{deepagents-cli}~\cite{deepagents}, a lightweight library providing tool-calling, skill loading, and response streaming via LangChain's message abstraction. Sage is LLM-agnostic: the active provider is specified in a user-level \texttt{config.toml} file. Supported backends include OpenAI (GPT-4o), Anthropic (Claude), Google (Gemini), and any OpenAI-compatible endpoint. The demonstrations in this paper use GLM-4.7, an open-weights model served on the NRP inference cluster, configured with temperature~$= 0$ for deterministic outputs.

Switching LLM provider requires only a \texttt{config.toml} change---no code modification. This design ensures that Sage is not tied to any specific model or cloud provider.

\subsection{The \texttt{\%\%ask} Magic Command}

The primary user interface is the \texttt{\%\%ask} cell magic. When a researcher writes a plain English prompt in a cell prefixed with \texttt{\%\%ask} and executes it, Sage: (1)~enriches the raw prompt with a system-level context injection; (2)~runs the \texttt{deepagents} agent asynchronously via \texttt{asyncio.get\_event\_loop().run\_until\_complete()}; (3)~streams tool call details and narration as the agent executes; and (4)~renders the agent's final report with inline maps, charts, and tables.

The prompt enrichment layer encodes several important design decisions. The \emph{one-map rule}---``your entire report must contain exactly one map tag; put every GeoJSON and WMS layer into that single tag as a comma-separated list''---prevents the agent from generating multiple disconnected maps when multiple spatial datasets are relevant. The \emph{output directory injection} tells the agent to save all files to \texttt{SAGE\_OUTPUT\_DIR}, preventing scattered outputs. The \emph{Python path injection} supplies \texttt{sys.executable} so the agent invokes the correct Python interpreter in \texttt{execute} tool calls.

A companion \texttt{\%reset} line magic clears both the output directory and conversation history, providing an explicit, user-controlled clean starting state.

\subsection{Cross-Cell Conversational Memory}

Unlike systems that treat each notebook cell as an isolated prompt, Sage maintains a Python list \texttt{SAGE\_MESSAGES} in the kernel's global namespace. Every \texttt{\%\%ask} cell appends the user's prompt and the agent's full response to this list. Subsequent \texttt{\%\%ask} prompts are prepended with the full conversation history before being sent to the agent.

This design has an important consequence: the agent in cell~$N$ has full access to everything discovered in cells 1 through~$N-1$ without the user needing to re-state context. In the earthquake/GNSS notebook, when the user asks ``Find GPS station datasets within 100 miles of the earthquake'' in cell~2, the agent resolves ``the earthquake'' from the history established in cell~1---the M5.7 Parker Butte event at a specific latitude/longitude---without any re-statement.

The memory model is deliberately stateless with respect to external storage: no SQLite checkpointer, no Redis, no database. Memory is exactly as portable as the Python process---it resets on kernel restart, persists across all \texttt{\%\%ask} cells in a session, and requires no infrastructure beyond the running kernel.

\subsection{Agent Execution and Tool Streaming}

The async function \texttt{\_run\_agent\_async} creates a \texttt{deepagents} agent with a \texttt{LocalShellBackend} (enabling real, unsandboxed code execution on the JupyterHub container), loads all installed skills from the skills directory, and streams the agent's response via \texttt{agent.astream()}.

As the agent streams, Sage separates two types of output: \emph{text narration} (the agent's reasoning and explanation, rendered progressively as Markdown) and \emph{tool calls} (intercepted before display and rendered as collapsible HTML panels showing the tool name, arguments, and result). Narration is buffered and flushed immediately before each tool call, producing the ``thinking out loud $\rightarrow$ action $\rightarrow$ result'' flow that mirrors how a scientist would document an investigation.

A deduplication mechanism silently discards replayed message chunks that the underlying agent framework may emit at the end of a streaming session, ensuring the final report is displayed exactly once.

\subsection{Output Directory Management and File Display}

Each notebook receives its own persistent output directory named \texttt{\_\{notebook\}\_sage\_/}, placed next to the notebook file and derived from the active kernel's session name. For a notebook named \texttt{flood\_impacts.ipynb}, all outputs are written to \texttt{\_flood\_impacts\_sage\_/}. This folder persists across kernel restarts and is cleared only on an explicit \texttt{\%reset}, making outputs durable across a working session while keeping each notebook's files self-contained.

After the agent completes, Sage scans the final report for standard Markdown image syntax \texttt{![caption](path)} and dispatches each file reference to the appropriate renderer: \texttt{.geojson} and \texttt{.wms.json} files become interactive Folium maps with a layer control; \texttt{.png} files become inline images. Multiple file paths in a single tag (comma-separated) are combined into a single multi-layer map, preserving the spatial relationship between datasets.

\subsection{Two-Tier Skill Architecture}

Domain knowledge in Sage is encoded as \emph{skills}---directories containing a single \texttt{SKILL.md} file written in natural language. A skill describes a data source's purpose, API patterns, output formats, and usage examples in free prose that the agent reads and reasons over at runtime. The agent selects which skills to invoke by matching the user's intent against the \texttt{description} field in each skill's YAML frontmatter; no vector index or embedding lookup is required.

This design has two important consequences. First, skills are hot-swappable: dropping a new \texttt{SKILL.md} into the skills directory makes it available in the next \texttt{\%\%ask} call with no code change or kernel restart. Second, because skills are plain text, they are legible, versionable, and composable by anyone---not just developers.

The skill architecture operates at two tiers:

\begin{itemize}
  \item \textbf{Bundled skills} are packaged in the Docker image at build time and cover the data sources used in the demonstration notebooks (see Section~4). They are available from the first \texttt{\%\%ask} cell of every session with no setup.
  \item \textbf{Runtime skills} are discovered and installed from the SkillsMP marketplace~\cite{skillsmp} during a live session via a natural language prompt. A single built-in \texttt{skillsmp} skill---the only system-level skill---gives every Sage instance access to 600,000+ community-contributed skills that can be browsed, inspected, and installed without leaving the notebook.
\end{itemize}

% ============================================================
\section{Demonstrations}

We present three complete demonstration notebooks, each targeting a distinct scientific use case and exercising different capabilities of the Sage system. All notebooks are available in the Sage GitHub repository~\cite{sage_github}. The demonstrations use data sources that are publicly accessible and not specific to any particular platform; the NDP catalog and NRP inference endpoint are used as convenient hosted services, but equivalents are freely substitutable.

\subsection{Earthquake and GNSS Analysis}

The \texttt{earthquake\_gnss} notebook demonstrates multi-step geophysical data integration across four \texttt{\%\%ask} cells, exercising cross-cell conversational memory. It uses two bundled skills: \texttt{usgs-earthquake-events} (for querying the USGS ComCat earthquake catalog) and \texttt{ndp-search} (for discovering GNSS station datasets in the NDP catalog).

\begin{enumerate}
  \item \textbf{Earthquake discovery}: ``Find earthquake events with magnitude $>$~5 in the United States between 2024-12-03 and 2024-12-15.'' Sage queries the USGS ComCat API, returns a structured table and GeoJSON map, and saves the events to the output directory.
  \item \textbf{GNSS station discovery}: ``Find GPS station datasets within 100 miles of the earthquake.'' Sage resolves the reference from \texttt{SAGE\_MESSAGES}, computes a bounding box around the M5.7 Parker Butte event, and searches the NDP OpenSearch catalog for GNSS datasets.
  \item \textbf{Data download}: Sage downloads raw GNSS time series data from the closest station, handling authentication, file transfer, and format detection autonomously.
  \item \textbf{Time series analysis}: Sage analyzes the downloaded GNSS displacement data around the earthquake date, identifies coseismic and postseismic signals in the north, east, and vertical displacement components, and produces annotated time series charts.
\end{enumerate}

This workflow illustrates Sage's key advantage over both manual coding and one-shot AI analysis: the chain of reasoning is explicit and inspectable at every step---from the USGS API query parameters to the NDP search bounding box to the rasterio sampling logic.

\subsection{Flood Impact Assessment: Kanawha River}

The \texttt{flood\_impacts} notebook demonstrates applied geospatial science, exercising WMS, WCS, and WFS data integration alongside the FEMA National Structure Inventory (NSI) to quantify human and economic impacts of a flood scenario on the Kanawha River basin. It uses four bundled skills: \texttt{kanawha-flood-depth} (USACE HEC-RAS 2D flood depth raster), \texttt{kanawha-nsi-impact} (NSI structure impact analysis), \texttt{kanawha-reach-impact} (reach-level flood extent), and \texttt{kanawha-cikr-impact} (critical infrastructure impact assessment).

\begin{enumerate}
  \item \textbf{Flood depth visualization}: Sage constructs a WMS URL pointing to USACE HEC-RAS 2D model results hosted on the SDSC GeoServer, overlays it on a Folium map, and renders an interactive flood extent visualization.
  \item \textbf{School impact analysis}: Sage fetches all 88,460 NSI 2022 structures via WFS, reprojects them to EPSG:5070, downloads the flood depth raster via WCS, and samples the raster at each structure location using rasterio, identifying flooded educational facilities.
  \item \textbf{Commercial building impact}: Sage filters for commercial occupancy types with structure value exceeding \$1 million, generating a summary table and spatial map of high-value assets within the flood footprint.
  \item \textbf{Elderly population at risk}: Sage cross-references NSI demographic attributes with flood depth to estimate the number of residents aged 65 and older within the flood footprint.
\end{enumerate}

The flood impact notebook demonstrates a complete, multi-source geospatial analysis pipeline---one that would typically require several hundred lines of GeoPandas, rasterio, and Folium code---expressed as four natural language prompts.

\subsection{On-the-Fly Skill Installation: Emergency Management}

The \texttt{skills\_manage} notebook demonstrates Sage's dynamic extensibility: discovering, installing, and immediately using a new domain skill within a single session, without restarting the kernel. It uses only the \texttt{skillsmp} skill---the single system-level skill bundled in every Sage image---which provides access to the SkillsMP marketplace of 600,000+ community-contributed skills.

\begin{enumerate}
  \item \textbf{Session reset}: \texttt{\%reset} clears the output directory and conversation history.
  \item \textbf{Marketplace search}: ``Search SkillsMP marketplace for an expert-level emergency management skill...'' Sage calls the SkillsMP AI search API and returns a ranked table of relevant skills.
  \item \textbf{Skill inspection}: ``Could you please show me the details of the skill \texttt{emergency-manager}?'' Sage retrieves and displays the skill's full documentation.
  \item \textbf{Skill installation}: ``Could you please install this skill for me?'' Sage downloads the \texttt{emergency-manager} skill to \texttt{\textasciitilde/.deepagents/agent/skills/}. It is immediately available; no kernel restart is required.
  \item \textbf{Immediate use}: ``We have a major hurricane approaching Florida next week. How do we coordinate the response?'' The newly installed skill provides a structured, expert-level response following NIMS/ICS principles.
\end{enumerate}

This demonstration illustrates a property of Sage with no analog in existing AI notebook systems: the system's capability envelope is user-extensible at runtime, in natural language, without code or configuration.

% ============================================================
\section{Design Principles and Discussion}

\subsection{Natural Language as the Universal Scientific Interface}

A recurring theme across all three demonstrations is that the user's scientific intent---expressed in plain English---is faithfully translated into a chain of technically sophisticated operations without requiring the user to write, review, or understand any code. A domain scientist who cannot write Python can conduct a complete multi-source geospatial analysis; a researcher unfamiliar with the USGS ComCat API can query it correctly on the first prompt.

This democratization is achievable precisely because Sage's skill architecture encodes domain knowledge in a form the agent can reason with. A skill is not a tool call specification; it is a domain expert's knowledge artifact---describing not just API syntax but scientific context, data quality considerations, appropriate output formats, and common analytical patterns.

\subsection{Transparency Through Traceability}

Sage's collapsible tool-call panels are not merely a user interface feature---they are a computational audit mechanism. Every \texttt{read\_file}, \texttt{execute}, and \texttt{http\_request} call is logged with its full parameters and output, inline in the notebook, adjacent to the natural language prompt that caused it.

The collapsible design is deliberate: tool logs are available for inspection without being intrusive during normal reading. A researcher presenting results to colleagues navigates the notebook at the prose-and-visualization level; a peer reviewer reproducing results can expand any tool call to inspect the exact query, code, or API parameters used.

\subsection{Reproducibility as an Architectural Property}

A central design decision in Sage is that reproducibility is not a feature added on top of the system---it is a property that emerges from the architecture. The combination of (1)~notebook-scoped persistent output directories, (2)~collapsible tool call logs embedded in the notebook, (3)~cross-cell conversational memory that records the full analytical chain, and (4)~standard \texttt{.ipynb} format means that a Sage notebook is inherently a reproducible scientific record: it contains the questions asked, the reasoning displayed, the tools invoked, and the outputs produced.

This contrasts sharply with conventional notebooks, where reproducibility requires explicit, manual effort: pinning library versions, documenting environment dependencies, exporting data snapshots, and writing narrative prose to explain the analytical chain.

\subsection{Human-in-the-Loop Scientific Agency}

Sage is intentionally positioned as a human-in-the-loop tool rather than a fully autonomous science agent. Each \texttt{\%\%ask} cell represents a deliberate scientific question posed by a human researcher. The agent executes one analytical step---it does not autonomously pursue a research agenda. The scientist retains control of the investigative trajectory at all times.

Fully autonomous science agents face deep challenges around hypothesis validity, methodological appropriateness, and the risk of generating plausible-but-incorrect results at scale. Sage sidesteps these challenges by keeping the human in the analytical loop: each cell's scope is bounded by the user's intent, and the tool logs make the agent's actions immediately auditable.

\subsection{Sage as the Claude Code of Science}

The most clarifying frame for Sage's contribution is the analogy to agentic software development tools. Claude Code and GitHub Copilot Workspace have demonstrated that AI agents can dramatically accelerate software development by automating the tedious, mechanical aspects of coding---API calls, boilerplate, debugging---while keeping the engineer in control of architecture and intent.

Sage applies the same agentic model to science, but changes the deliverable. The scientist's deliverable is not working code but a reproducible scientific record---a notebook that documents the question, the analytical chain, and the result in a form that a colleague can inspect, validate, and re-run. Sage automates the tedious aspects of scientific data work---API wrangling, format conversion, spatial projection, visualization---while keeping the scientist in control of the scientific questions and their interpretation.

% ============================================================
\section{Limitations and Future Work}

\subsection{Current Limitations}

\begin{itemize}
  \item \textbf{Setup configuration}: While Sage runs on any Jupyter environment, configuring a new LLM provider requires editing \texttt{config.toml} and supplying an API key. A one-click setup wizard or auto-detection of available providers is not yet implemented.
  \item \textbf{LLM non-determinism}: While Sage's outputs are structurally reproducible, the specific code the agent generates and the natural language narrative may vary between runs. Temperature~$= 0$ mitigates but does not eliminate this.
  \item \textbf{Skill quality}: The quality of Sage's domain reasoning is bounded by the quality of installed skills. Community-contributed marketplace skills vary in documentation depth, accuracy, and currency.
  \item \textbf{No formal evaluation}: This paper presents system demonstrations rather than a controlled user study or benchmarked evaluation.
\end{itemize}

\subsection{Future Work}

\begin{itemize}
  \item \textbf{Multi-user collaborative notebooks}: Extending Sage's conversational memory to shared, multi-user sessions for collaborative scientific analysis.
  \item \textbf{Skill authoring from within Sage}: Allowing users to define and publish new skills directly from \texttt{\%\%ask} prompts.
  \item \textbf{Automated correctness evaluation}: Connecting Sage to domain-specific validation oracles for benchmarking.
  \item \textbf{Data provenance standards integration}: Mapping tool call logs to W3C PROV-O or schema.org Dataset standards for machine-readable audit trails.
  \item \textbf{One-click deployment}: Simplifying Sage setup for users unfamiliar with Docker or \texttt{config.toml}, including automatic LLM provider detection and guided first-run configuration.
  \item \textbf{Formal user study}: Comparing Sage to conventional notebook workflows on representative scientific tasks.
\end{itemize}

% ============================================================
\section{Conclusion}

We have presented Sage, a natural language science agent that makes AI-driven data exploration traceable, reproducible, and shareable by architectural design. By embedding a persistent conversational memory model, a collapsible tool-call audit log, and a notebook-scoped output management system in a standard JupyterHub environment, Sage transforms the Jupyter notebook from a display medium into an auditable scientific record.

The three demonstrations presented here---geophysical event analysis, flood impact assessment, and on-the-fly skill installation---illustrate that Sage is not merely a convenience layer over existing tools, but a new approach to scientific computing: one in which the unit of work is a natural language question, the unit of output is a reproducible notebook cell, and the unit of extension is a community-contributed skill.

Reproducibility is not a feature to be added to a science agent---it is an architectural property that must be designed in from the start. When an AI agent's execution is embedded in a notebook that records every tool call, preserves every output, and maintains the full conversational context of a scientific investigation, reproducibility becomes the default rather than the exception.

We believe Sage represents a significant step toward a future in which the barrier between a scientific question and a reproducible, shareable, verifiable answer is as low as typing a sentence in plain English.

% ============================================================
\section*{Acknowledgments}

\noindent [To be completed. Suggested: National Data Platform (NDP), National Research Platform (NRP), NSF awards supporting NDP/NRP infrastructure, WENOKN project, San Diego Supercomputer Center (SDSC), SkillsMP.]

% ============================================================
\bibliographystyle{plain}
\begin{thebibliography}{99}

\bibitem{jupyter_ai}
Jupyter AI. \textit{Official AI Extension for JupyterLab}. \url{https://github.com/jupyterlab/jupyter-ai}

\bibitem{deepagents}
deepagents-cli. \textit{A Lightweight Agent Framework with Tool Calling and Skill Loading}. \url{https://pypi.org/project/deepagents-cli/}

\bibitem{deepnote}
Deepnote. \textit{AI-powered Data Notebook}. \url{https://deepnote.com}

\bibitem{chemcrow}
A. M. Bran et al., ``ChemCrow: Augmenting Large-Language Models with Chemistry Tools,'' \textit{Nature Machine Intelligence}, 2024.

\bibitem{skillsmp}
SkillsMP. \textit{The AI Agent Skills Marketplace}. \url{https://skillsmp.com}

\bibitem{sage_github}
Sage GitHub Repository. \url{https://github.com/klinucsd/sage}

\bibitem{ndp}
National Data Platform. \textit{NDP: A Federated and Extensible Library of Datasets for Science}. \url{https://nationaldataplatform.org}

\bibitem{nrp}
National Research Platform. \textit{NRP Nautilus Hypercluster}. \url{https://nrp-nautilus.io}

\bibitem{hec_ras}
U.S. Army Corps of Engineers. \textit{HEC-RAS 2D: Two-Dimensional River Analysis System}. \url{https://www.hec.usace.army.mil/software/hec-ras/}

\bibitem{nsi}
FEMA. \textit{National Structure Inventory (NSI) 2022}. \url{https://www.fema.gov/flood-maps/products-tools/national-structure-inventory}

\bibitem{usgs_comcat}
USGS. \textit{ComCat: The USGS Comprehensive Earthquake Catalog}. \url{https://earthquake.usgs.gov/data/comcat/}

\bibitem{perkel2018}
J. M. Perkel, ``Why Jupyter is data scientists' computational notebook of choice,'' \textit{Nature}, vol.~563, no.~7729, pp.~145--146, 2018.

\bibitem{pimentel2019}
J. F. Pimentel et al., ``A Large-Scale Study About Quality and Reproducibility of Jupyter Notebooks,'' in \textit{Proc. IEEE/ACM MSR}, 2019.

\bibitem{hou2024}
X. Hou et al., ``Large Language Models for Software Engineering: A Systematic Literature Review,'' \textit{ACM Trans. Softw. Eng. Methodol.}, 2024.

\bibitem{w3c_provo}
W3C. \textit{PROV-O: The PROV Ontology}. \url{https://www.w3.org/TR/prov-o/}

\end{thebibliography}

\end{document}
